профиль 2020 (алюминий) & 15 & 500 мм ×2, 485 мм ×2, 460 мм ×3, 252 мм ×2, 82 мм ×2, 50 мм ×4\\
установочный винт М3 с плоским концом & 103 & ×6: 2, ×8: 2, ×20: 1, ×25: 98\\
стальной пруток 5 мм (планки рамок) & 8 & 228 мм ×2, 238 мм ×2, 276 мм ×2, 306 мм ×2\\
линейная направляющая MGN12, 300 мм & 2 & опоры ползуна (U-рамы)\\
каретка MGN12H & 4 & \\
линейная направляющая MGN7, 75 мм & 1 & каретка ленты поворота\\
каретка MGN7H & 1 & \\
калиброванный пруток 8 мм, 460 мм & 2 & направляющие каретки\\
линейный подшипник LM8UU & 4 & каретка\\
подшипник 608 & 2 & вал челнока\\
заводная пружина 32×0,8 мм, около 6,3 м & 1 & граммофонного типа, с концевыми петлями\\
вал 10 мм, 100 мм & 1 & ось барабана пружины\\
вал 8 мм, 380 мм & 1 & цикловой вал\\
бронзовая втулка 6 мм & 1 & поводок челнока\\
стальная лента 5×0,1 мм & 3 & лента поворота, около 0,15 м\\
пруток 3 мм, 40 мм & 7 & стержни выбора\\
пруток 3 мм, 80 мм & 6 & толкатели кулачков\\
пруток 4 мм, 200 мм & 1 & промежуточный вал\\
плита Д16Т 326×188×6 & 1 & основание печатающего узла, сверловка\\
вал 6 мм с резиновой трубкой & 1 & подающий валик\\
прижимной ролик 7 мм & 2 & \\
труба 30 мм, 230 мм & 1 & кассета бумаги\\
катушка ленты 13 мм & 2 & красящая лента\\
пруток 3--6 мм, короткие оси & 11 & оси секторов, молотка, спуска, передачи, поводок\\
вал 8 мм & 1 & вал челнока, 160 мм\\
rod2x10 & 3 & \\
rod2x20 & 1 & \\
rod3x12 & 2 & \\
брус гребёнок: квадрат 6×6 мм, на двух стойках; гребёнки клавиш и пробела на лапках, на винтах М3 (ред. 0.14) & 1 & \\
ведущая планка 3 мм, ход 12 мм & 2 & \\
вибратор ленты: вилка из жести 0,5 мм на ползуне; ход 10 мм & 1 & \\
войлочное кольцо 13×5×0,6 мм, наклеено на шайбу & 1 & \\
грузик регулятора: латунный шарик Ø4 мм на пружине & 3 & \\
зажим троса W на площадке тяги Y (винт нуля) & 1 & \\
запорная гребёнка 2 мм: прорези под рычаги, сдвиг на полшага на время цикла & 1 & \\
клин фиксатора ряда, 60°, ход 5 мм & 1 & \\
кулачок: сталь 45 5 мм, закалка; ролик закалённый (ред. 0.14) & 2 & \\
маховик регулятора: стальной диск 31×3,2 мм на оси регулятора (ред. 0.14) & 1 & \\
наконечник оболочки троса Ø3,5 & 4 & \\
наконечник троса освобождения каретки на хвосте собачки & 1 & \\
направляющая бумаги -- полоса 0,5 мм & 1 & \\
направляющие 3 мм & 12 & \\
направляющие ползунка 2 мм & 2 & \\
неподвижная муфта: латунь, штифт 1 мм в оси & 1 & \\
обратный кулачок: сталь 45 5 мм, закалка; профиль -- основной через коромысло, в выстое отступ 0,5 мм (ред. 0.14) & 2 & \\
основание: фанера 4 мм под рамой, на винтах М5 с Т-гайками в нижние пазы профиля; детали, стоящие на основании, привёрнуты снизу винтами М3 во вставки (ред. 0.14) & 1 & \\
ось 4 мм, запрессована в молоток (ред. 0.14) & 1 & \\
ось 4 мм, запрессована в правую щёку двигателя & 2 & \\
ось защёлки цикла (ред. 0.14) & 1 & \\
ось ролика (ред. 0.14) & 3 & \\
ось шептала (ред. 0.14) & 1 & \\
палец выходного ползуна & 4 & \\
палец защёлки 3 мм с пружиной вниз & 7 & \\
плоская пружина 0,1×2 мм между муфтами & 3 & \\
ползун опоры: смещение нуля от ведущей планки & 3 & \\
ползун разряда: полоса 4×3 мм & 7 & \\
ползун узла коромысла: полоса 4×3 мм & 5 & \\
полоса 4×6 мм, 222 мм, с бобышками под пружины (запрессованные штифты) & 1 & \\
пружина кручения между рычагом взвода и молотком (ред. 0.14) & 1 & \\
пружина растяжения возврата узла, около 1 Н & 1 & \\
пружина растяжения клина ряда, 0,5 Н (ред. 0.14) & 1 & \\
пружина сжатия 3,5×0,3 мм, 0,34 Н в покое, 0.01 Н/мм & 46 & \\
пружина сжатия 5,2×0,5 мм на направляющей & 8 & \\
пружина сжатия: поджатие 1,5 Н, до 2,3 Н & 1 & \\
пруток 4 мм, 315 мм, с канавками под стопорные кольца & 1 & \\
ролик на клине ряда & 1 & \\
ролик обратной ветви ленты (ред. 0.14) & 1 & \\
скользящая муфта: латунь, торец по войлоку & 1 & \\
стопорное кольцо DIN 6799 3,2 на оси рычагов (ред. 0.14) & 2 & \\
стопорное кольцо на оси запора чтения (ред. 0.14) & 1 & \\
стопорное кольцо на прутке каретки щупа (ред. 0.14) & 1 & \\
тяга 1,5 мм с шаровыми наконечниками и пружиной, упорная втулка & 2 & \\
тяга 2 мм к упору спуска, углепластиковый пруток (ред. 0.14) & 1 & \\
тяга 2 мм, 430 мм, углепластиковый пруток (работает на растяжение; ред. 0.14: вместо стали -- щупу W вчетверо легче) & 1 & \\
тяга A, задняя часть (паз под палец повторителя) & 3 & \\
тяга вибратора: кулачок ленты поднимает её на 10 мм & 1 & \\
тяга запорной гребёнки, полоса 3×2,5 мм & 1 & \\
тяга от качалки к ведущей планке, палец в пазу качалки & 1 & \\
тяга отвода щупов (полоса 1,5 мм) & 1 & \\
тяга спуска цикла к защёлке & 1 & \\
указатель точки печати: проволока 0,8 мм с красной бусиной 1,6 мм на головке чтения & 1 & \\
упор нуля M: стальной штифт 3 мм в мостике над линейкой полей (ред. 0.14) & 1 & \\
штифт 2×5 мм (гребёнка на кольце) & 22 & \\
штырь запора удара Ø3×18 с головкой Ø7 и резиновой втулкой; входит на 6,2 мм перед плечом молотка & 1 & \\
штырь запора чтения Ø3×14 с головкой Ø7; входит на 5 мм перед нижним плечом коромысла & 1 & \\
щуп упоров: полоса 4×4 мм на плите узла & 1 & \\
щуп: полоса 1,6×2×22 мм, лапка с зажимом троса; на носу эластомерный наконечник 0,5 мм (ред. 0.14, ступень 2б: отскок от зубца ≤ 0,3 мм) & 3 & \\
